\section{Lifecycle Pengembangan Game}

\begin{learningoutcomes}
\item Memahami fase-fase dalam lifecycle pengembangan game
\item Mengetahui aktivitas utama setiap fase
\item Memahami pentingnya planning dan documentation
\item Mampu mengidentifikasi milestone dalam development
\end{learningoutcomes}

\subsection{Pengertian Lifecycle Pengembangan Game}

\begin{definisi}[Game Development Lifecycle]
Lifecycle pengembangan game adalah proses sistematis dari konsep awal hingga rilis dan maintenance game, yang terdiri dari beberapa fase utama \cite{Gregory2018}.
\end{definisi}

\subsection{Fase-Fase Pengembangan Game}

\subsubsection{1. Pre-Production}

\textbf{Tujuan}: Membangun fondasi dan visi game

\textbf{Aktivitas Utama}: Brainstorming konsep, market research, prototyping core mechanics, menyusun GDD, budgeting dan resource planning, team assembly.

\textbf{Output}: GDD lengkap, prototype yang dapat dimainkan, art bible dan style guide, technical design document, project plan dan timeline.

\subsubsection{2. Production}

\textbf{Tujuan}: Membangun konten dan fitur game

\textbf{Aktivitas Utama}: Asset creation (art, audio, models), programming dan implementation, level design, integration dan testing, iterasi berdasarkan feedback.

\textbf{Fase dalam Production}:
\begin{enumerate}
\item \textbf{Alpha}: Core features implemented, masih banyak bug
\item \textbf{Beta}: Feature-complete, focus pada bug fixing dan polish
\item \textbf{Gold Master}: Final version siap untuk rilis
\end{enumerate}

\subsubsection{3. Post-Production}

\textbf{Tujuan}: Rilis, marketing, dan maintenance

\textbf{Aktivitas Utama}: Final testing dan QA, marketing dan promotion, launch preparation, distribution, post-launch support, updates dan patches, DLC development.

\subsection{Metodologi Pengembangan}

\textbf{Waterfall}: Pendekatan linear dengan fase berurutan. Keuntungan: terstruktur. Kekurangan: kurang fleksibel.

\textbf{Agile/Scrum}: Pendekatan iteratif dengan sprint pendek. Keuntungan: fleksibel, adaptif. Kekurangan: memerlukan komunikasi intensif.

\textbf{Hybrid}: Kombinasi waterfall untuk planning dan agile untuk execution.

\subsection{Milestone Penting}

Concept Approval, Prototype Complete, Pre-Alpha, Alpha, Beta, Gold, Launch.

\subsection{Best Practices}

Start Small, Prototype Early, Iterate Often, Document Everything, Test Continuously.
